\chapter{2005年江西卷（理）}
\section{单选题}
\begin{problem}
设集合 \(I=\{x\mid |x|<3,\ x\in\Z\}\)，\(A=\{1,2\}\)，\(B=\{-2,-1,2\}\)，则 \(A\cup(\complement_I B)=\)（\quad）

\choices
    {\(\{1\}\)}
    {\(\{1,2\}\)}
    {\(\{2\}\)}
    {\(\{0,1,2\}\)}
\end{problem}

\begin{answer}D\end{answer}

\begin{problem}
设复数：\(z_1=1+\i\)，\(z_2=x+2\i\)（\(x\in\R\)）.若 \(z_1z_2\) 为实数，则 \(x=\)（\quad）

\choices
    {\(-2\)}
    {\(-1\)}
    {\(1\)}
    {\(2\)}
\end{problem}

\begin{answer}A\end{answer}

\begin{problem}
“\(a=b\)”是“直线 \(y=x+2\) 与圆 \((x-a)^2+(y+b)^2=2\) 相切”的（\quad）

\choices
    {充分不必要条件}
    {必要不充分条件}
    {充分必要条件}
    {既不充分又不必要条件}
\end{problem}

\begin{answer}D\end{answer}

\begin{problem}
\((\sqrt{x}+\sqrt[3]{x})^{12}\) 的展开式中，含 \(x\) 的正整数次幂的项共有（\quad）

\choices
    {\(4\) 项}
    {\(3\) 项}
    {\(2\) 项}
    {\(1\) 项}
\end{problem}

\begin{answer}B\end{answer}

\begin{problem}
设函数 \(f(x)=\sin 3x+|\sin 3x|\)，则 \(f(x)\) 为（\quad）

\choices
    {周期函数，最小正周期为 \(\frac{\pi}{3}\)}
    {周期函数，最小正周期为 \(\frac{2\pi}{3}\)}
    {周期函数，最小正周期为 \(2\pi\)}
    {非周期函数}
\end{problem}

\begin{answer}B\end{answer}

\begin{problem}
已知向量 \(\bs{a}=(1,2)\)，\(\bs{b}=(-2,-4)\)，\(|\bs{c}|=\sqrt{5}\)，若 \((\bs{a}+\bs{b})\cdot\bs{c}=\frac{5}{2}\)，则 \(\bs{a}\) 与 \(\bs{c}\) 的夹角为（\quad）

\choices
    {\(30^{\circ}\)}
    {\(60^{\circ}\)}
    {\(120^{\circ}\)}
    {\(150^{\circ}\)}
\end{problem}

\begin{answer}C\end{answer}

\begin{problem}
已知函数 \(y=xf'(x)\) 的图象如图所示（其中 \(f'(x)\) 是函数 \(f(x)\) 的导函数），下面四个图象中，\(y=f(x)\) 的图象大致是（\quad）

\bitmapfigure[max width=0.42\linewidth]{img/2005/jiangxi_science/q7_fig1.png}

\choices
    {\choicebitmap[width=0.15\paperwidth]{img/2005/jiangxi_science/q7_optA_clean.png}}
    {\choicebitmap[width=0.15\paperwidth]{img/2005/jiangxi_science/q7_optB_clean.png}}
    {\choicebitmap[width=0.15\paperwidth]{img/2005/jiangxi_science/q7_optC_clean.png}}
    {\choicebitmap[width=0.15\paperwidth]{img/2005/jiangxi_science/q7_optD_clean.png}}
\end{problem}

\begin{answer}C\end{answer}

\begin{problem}
若 \(\displaystyle\lim_{x\to 1}\frac{f(x-1)}{x-1}=1\)，则 \(\displaystyle\lim_{x\to 1}\frac{x-1}{f(2-2x)}=\)（\quad）

\choices
    {\(-1\)}
    {\(1\)}
    {\(-\frac{1}{2}\)}
    {\(\frac{1}{2}\)}
\end{problem}

\begin{answer}C\end{answer}

\begin{problem}
矩形 \(ABCD\) 中，\(AB=4\)，\(BC=3\)，沿 \(AC\) 将矩形 \(ABCD\) 折成一个直二面角 \(B-AC-D\)，则四面体 \(ABCD\) 的外接球的体积为（\quad）

\choices
    {\(\frac{125}{12}\pi\)}
    {\(\frac{125}{9}\pi\)}
    {\(\frac{125}{6}\pi\)}
    {\(\frac{125}{3}\pi\)}
\end{problem}

\begin{answer}C\end{answer}

\begin{problem}
已知实数 \(a\)，\(b\) 满足等式 \(\left(\frac{1}{2}\right)^a=\left(\frac{1}{3}\right)^b\)，下列五个关系式：

\begin{circlelist}
\item \(0<b<a\)；
\item \(a<b<0\)；
\item \(0<a<b\)；
\item \(b<a<0\)；
\item \(a=b\).
\end{circlelist}

其中不可能成立的关系式有（\quad）

\choices
    {\(1\) 个}
    {\(2\) 个}
    {\(3\) 个}
    {\(4\) 个}
\end{problem}

\begin{answer}B\end{answer}

\begin{problem}
在 \(\triangle OAB\) 中，\(O\) 为坐标原点，\(A(1,\cos\theta)\)，\(B(\sin\theta,1)\)，\(\theta\in\left(0,\frac{\pi}{2}\right]\)，则 \(\triangle OAB\) 的面积达到最大值时，\(\theta=\)（\quad）

\choices
    {\(\frac{\pi}{6}\)}
    {\(\frac{\pi}{4}\)}
    {\(\frac{\pi}{3}\)}
    {\(\frac{\pi}{2}\)}
\end{problem}

\begin{answer}D\end{answer}

\begin{problem}
将 \(1\)，\(2\)，\(\cdots\)，\(9\) 这 \(9\) 个数平均分成三组，则每组的三个数都成等差数列的概率为（\quad）

\choices
    {\(\frac{1}{56}\)}
    {\(\frac{1}{70}\)}
    {\(\frac{1}{336}\)}
    {\(\frac{1}{420}\)}
\end{problem}

\begin{answer}A\end{answer}

\section{填空题}
\begin{problem}
若函数 \(f(x)=\log_{a}\left(x+\sqrt{x^{2}+2a^{2}}\right)\) 是奇函数，则 \(a=\)\fillinblank{}.
\end{problem}

\begin{answer}\(\frac{\sqrt{2}}{2}\)\end{answer}

\begin{problem}
设实数 \(x\)，\(y\) 满足 \(\begin{cases}
x-y-2\leqslant 0,\\
 x+2y-4\geqslant 0,\\
2y-3\leqslant 0
\end{cases}\)，则 \(\frac{y}{x}\) 的最大值是\fillinblank{}.
\end{problem}

\begin{answer}\(\frac{3}{2}\)\end{answer}

\begin{problem}
如图，在直三棱柱 \(ABC-A_{1}B_{1}C_{1}\) 中，\(AB=BC=\sqrt{2}\)，\(BB_{1}=2\)，\(\angle ABC=90^{\circ}\)，\(E\)，\(F\) 分别为 \(AA_{1}\)，\(C_{1}B_{1}\) 的中点，沿棱柱的表面从 \(E\) 到 \(F\) 两点的最短路径的长度为\fillinblank{}.

\bitmapfigure[max width=0.42\linewidth]{img/2005/jiangxi_science/q15_fig1.png}
\end{problem}

\begin{answer}\(\frac{3\sqrt{2}}{2}\)\end{answer}

\begin{problem}
以下四个关于圆锥曲线的命题中：

\begin{circlelist}
\item 设 \(A\)，\(B\) 为两个定点，\(k\) 为非零常数，\(\left|\overrightarrow{PA}\right|-\left|\overrightarrow{PB}\right|=k\)，则动点 \(P\) 的轨迹为双曲线；
\item 设圆 \(C\) 上一定点 \(A\) 作圆的动点弦 \(AB\)，\(O\) 为坐标原点，若 \(\overrightarrow{OP}=\frac{1}{2}\left(\overrightarrow{OA}+\overrightarrow{OB}\right)\)，则动点 \(P\) 的轨迹为椭圆；
\item 方程 \(2x^{2}-5x+2=0\) 的两根可分别作为椭圆和双曲线的离心率；
\item 双曲线 \(\frac{x^{2}}{25}-\frac{y^{2}}{9}=1\) 与椭圆 \(\frac{x^{2}}{35}+y^{2}=1\) 有相同的焦点.
\end{circlelist}

其中真命题的序号为\fillinblank{}.（写出所有真命题的序号）
\end{problem}

\begin{answer}\(\circled{3}\circled{4}\)\end{answer}

\section{解答题}
\begin{problem}
已知函数 \(f(x)=\frac{x^{2}}{ax+b}\)（\(a\)，\(b\) 为常数），且方程 \(f(x)-x+12=0\) 有两个实根为 \(x_1=3\)，\(x_2=4\).

\begin{enumerate}
\item 求函数 \(f(x)\) 的解析式；
\item 设 \(k>1\)，解关于 \(x\) 的不等式：\(f(x)<\frac{(k+1)x-k}{2-x}\).
\end{enumerate}
\end{problem}

\begin{solution}
\begin{enumerate}
\item 将 \(x=3\)，\(x=4\) 分别代入方程，得
\[
\begin{cases}
\dfrac{9}{3a+b}-3+12=0,\\
\dfrac{16}{4a+b}-4+12=0.
\end{cases}
\]
因此
\[
\begin{cases}
3a+b=-1,\\
4a+b=-2,
\end{cases}
\]
解得 \(a=-1\)，\(b=2\)，所以 \(f(x)=\dfrac{x^{2}}{2-x}\)，其中 \(x\ne2\). 直接计算得
\[
f(x)-x+12=\frac{x^{2}}{2-x}-x+12=\frac{2(x-3)(x-4)}{2-x},
\]
故原方程的两个实根确为 \(3\)，\(4\).

\item 由定义域知 \(x\ne2\)，不等式等价于
\[
\frac{x^{2}}{2-x}<\frac{(k+1)x-k}{2-x}
\quad\Longleftrightarrow\quad
\frac{x^{2}-(k+1)x+k}{2-x}<0
\quad\Longleftrightarrow\quad
\frac{(x-1)(x-k)}{2-x}<0.
\]
当 \(1<k\leqslant2\) 时，关键点的顺序为 \(1<k\leqslant2\)，作符号判断得 \(x\in(1,k)\cup(2,+\infty)\). 当 \(k>2\) 时，关键点的顺序为 \(1<2<k\)，作符号判断得 \(x\in(1,2)\cup(k,+\infty)\). 因此不等式的解集为
\[
x\in
\begin{cases}
(1,k)\cup(2,+\infty),&1<k\leqslant2,\\
(1,2)\cup(k,+\infty),&k>2.
\end{cases}
\]
\end{enumerate}
\end{solution}

\begin{problem}
已知向量 \(\bs{a}=\left(2\cos\frac{x}{2},\tan\left(\frac{x}{2}+\frac{\pi}{4}\right)\right)\)，\(\bs{b}=\left(\sqrt{2}\sin\left(\frac{x}{2}+\frac{\pi}{4}\right),\tan\left(\frac{x}{2}-\frac{\pi}{4}\right)\right)\)，令 \(f(x)=\bs{a}\cdot\bs{b}\). 是否存在实数 \(x\in[0,\pi]\)，使 \(f(x)+f'(x)=0\)（其中 \(f'(x)\) 是 \(f(x)\) 的导函数）？若存在，则求出 \(x\) 的值；若不存在，则证明之.
\end{problem}

\begin{solution}
令 \(t=\dfrac{x}{2}\). 由于 \(\tan\left(t+\dfrac{\pi}{4}\right)\) 有意义，必须有 \(t\ne\dfrac{\pi}{4}\)，即 \(x\ne\dfrac{\pi}{2}\). 在定义域内，利用
\[
\tan\left(t+\frac{\pi}{4}\right)\tan\left(t-\frac{\pi}{4}\right)=-1
\]
及 \(\sqrt{2}\sin\left(t+\dfrac{\pi}{4}\right)=\sin t+\cos t\)，得
\[
f(x)=2\cos t\left(\sin t+\cos t\right)-1=\sin x+\cos x.
\]
因此 \(f'(x)=\cos x-\sin x\)，从而
\[
f(x)+f'(x)=2\cos x.
\]
若 \(f(x)+f'(x)=0\)，则 \(\cos x=0\)，在 \([0,\pi]\) 内只能有 \(x=\dfrac{\pi}{2}\). 但该值不在函数定义域内，故不存在满足条件的实数 \(x\).
\end{solution}

\begin{problem}
\(A\)，\(B\) 两位同学各有五张卡片，现以投掷均匀硬币的形式进行游戏，当出现正面朝上时 \(A\) 赢得 \(B\) 一张卡片，否则 \(B\) 赢得 \(A\) 一张卡片. 规定掷硬币的次数达 \(9\) 次时，或在此前某人已赢得所有卡片时游戏终止. 设 \(\xi\) 表示游戏终止时掷硬币的次数.

\begin{enumerate}
\item 求 \(\xi\) 的取值范围；
\item 求 \(\xi\) 的数学期望 \(E(\xi)\).
\end{enumerate}
\end{problem}

\begin{solution}
设前 \(n\) 次投掷中正面和反面出现的次数分别为 \(m\)，\(r\). 某人赢得全部五张卡片时，正面与反面的次数之差的绝对值为 \(5\)，所以在游戏未因达到九次而终止时，必须有 \(n\geqslant5\)，且 \(n\) 与 \(5\) 同奇偶. 结合最多投掷九次，得到
\[
\xi\in\{5,7,9\}.
\]

当 \(\xi=5\) 时，前五次必须全为正面或全为反面，共有 \(2\) 个等可能序列，因此
\[
P(\xi=5)=2\left(\frac12\right)^5=\frac1{16}.
\]

当 \(\xi=7\) 时，以正面为例，七次中应有六次正面和一次反面，且反面必须出现在前五次中的一个位置，否则前五次正面会使 \(A\) 提前赢得全部卡片. 反面的位置有 \(5\) 种；以反面为主的情形同理，故共有 \(10\) 个等可能序列，从而
\[
P(\xi=7)=10\left(\frac12\right)^7=\frac5{64}.
\]

其余情形均在第九次终止，所以
\[
P(\xi=9)=1-P(\xi=5)-P(\xi=7)=1-\frac1{16}-\frac5{64}=\frac{55}{64}.
\]
因此
\[
E(\xi)=5P(\xi=5)+7P(\xi=7)+9P(\xi=9)=\frac{275}{32}.
\]
\end{solution}

\begin{problem}
如图，在长方体 \(ABCD-A_{1}B_{1}C_{1}D_{1}\) 中，\(AD=AA_{1}=1\)，\(AB=2\)，点 \(E\) 在棱 \(AB\) 上移动.

\begin{enumerate}
\item 证明：\(D_{1}E\perp A_{1}D\)；
\item 当 \(E\) 为 \(AB\) 的中点时，求点 \(E\) 到面 \(ACD_{1}\) 的距离；
\item \(AE\) 等于何值时，二面角 \(D_{1}-EC-D\) 的大小为 \(\frac{\pi}{4}\).
\end{enumerate}

\bitmapfigure[max width=0.56\linewidth]{img/2005/jiangxi_science/q20_fig1.png}
\end{problem}

\begin{solution}
建立空间直角坐标系，取 \(A=(0,0,0)\)，\(B=(2,0,0)\)，\(D=(0,1,0)\)，\(A_{1}=(0,0,1)\). 于是 \(C=(2,1,0)\)，\(D_{1}=(0,1,1)\). 设 \(AE=t\)，则 \(E=(t,0,0)\)，且 \(0\leqslant t\leqslant2\).

\begin{enumerate}
\item \(\overrightarrow{D_{1}E}=(t,-1,-1)\)，\(\overrightarrow{A_{1}D}=(0,1,-1)\)，所以
\[
\overrightarrow{D_{1}E}\cdot\overrightarrow{A_{1}D}=0-1+1=0.
\]
因此 \(D_{1}E\perp A_{1}D\).

\item 平面 \(ACD_{1}\) 的一个法向量为
\[
\overrightarrow{AC}\times\overrightarrow{AD_{1}}=(1,-2,2)
\]
故其方程为 \(x-2y+2z=0\). 当 \(E\) 为 \(AB\) 的中点时，\(E=(1,0,0)\)，所以点 \(E\) 到平面 \(ACD_{1}\) 的距离为
\[
\frac{|1-2\cdot0+2\cdot0|}{\sqrt{1^2+(-2)^2+2^2}}=\frac13.
\]

\item 平面 \(ECD\) 是底面，法向量可取 \(\bs n_0=(0,0,1)\). 平面 \(D_{1}EC\) 的两个方向向量为 \(\overrightarrow{EC}=(2-t,1,0)\)，\(\overrightarrow{ED_{1}}=(-t,1,1)\)，故可取法向量
\[
\bs n_1=\overrightarrow{EC}\times\overrightarrow{ED_{1}}=(1,t-2,2).
\]
设二面角 \(D_{1}-EC-D\) 为 \(\varphi\)，则
\[
\cos\varphi=\frac{|\bs n_0\cdot\bs n_1|}{|\bs n_0||\bs n_1|}=\frac{2}{\sqrt{5+(t-2)^2}}.
\]
当 \(\varphi=\dfrac{\pi}{4}\) 时，
\[
\frac{2}{\sqrt{5+(t-2)^2}}=\frac{\sqrt2}{2}
\quad\Longrightarrow\quad
(t-2)^2=3.
\]
由于 \(0\leqslant t\leqslant2\)，得 \(t=2-\sqrt3\). 所以 \(AE=2-\sqrt3\).
\end{enumerate}
\end{solution}

\begin{problem}
已知数列 \(\{a_n\}\) 的各项都是正数，且满足：\(a_{0}=1\)，\(\ a_{n+1}=\frac{1}{2}a_n(4-a_n)\)，\(n\in\mathbf{N}\).

\begin{enumerate}
\item 证明：\(a_n<a_{n+1}<2\)，\(\ n\in\mathbf{N}\)；
\item 求数列 \(\{a_n\}\) 的通项公式 \(a_n\).
\end{enumerate}
\end{problem}

\begin{solution}
\begin{enumerate}
\item 因为 \(0<a_0=1<2\)，设对某个 \(n\) 有 \(0<a_n<2\). 则
\[
a_{n+1}=\frac12a_n(4-a_n)>0,
\qquad
2-a_{n+1}=2-\frac12a_n(4-a_n)=\frac12(2-a_n)^2>0,
\]
所以 \(0<a_{n+1}<2\). 又
\[
a_{n+1}-a_n=\frac12a_n(4-a_n)-a_n=\frac12a_n(2-a_n)>0,
\]
故 \(a_n<a_{n+1}<2\). 由数学归纳法，对一切 \(n\in\mathbf N\) 命题成立.

\item 令 \(b_n=1-\dfrac{a_n}{2}\)，则
\[
b_{n+1}=1-\frac{a_{n+1}}2=1-a_n+\frac{a_n^2}{4}=\left(1-\frac{a_n}{2}\right)^2=b_n^2.
\]
又 \(b_0=\dfrac12\)，所以 \(b_n=\left(\dfrac12\right)^{2^n}\). 因此
\[
a_n=2(1-b_n)=2\left(1-\left(\frac12\right)^{2^n}\right).
\]
\end{enumerate}
\end{solution}

\begin{problem}
如图，设抛物线 \(C:y=x^{2}\) 的焦点为 \(F\)，动点 \(P\) 在直线 \(l:x-y-2=0\) 上运动，过 \(P\) 作抛物线 \(C\) 的两条切线 \(PA\)，\(PB\)，且与抛物线 \(C\) 分别相切于 \(A\)，\(B\) 两点.

\begin{enumerate}
\item 求 \(\triangle APB\) 的重心 \(G\) 的轨迹方程；
\item 证明 \(\angle PFA=\angle PFB\).
\end{enumerate}

\bitmapfigure[max width=0.58\linewidth]{img/2005/jiangxi_science/q22_fig1.png}
\end{problem}

\begin{solution}
设 \(P=(p,p-2)\)，并设切点 \(A=(u,u^2)\)，\(B=(v,v^2)\). 抛物线 \(y=x^2\) 在参数点 \((r,r^2)\) 处的切线为 \(y=2rx-r^2\)，所以 \(u\)，\(v\) 是方程
\[
r^2-2pr+p-2=0
\]
的两个根，从而
\[
u+v=2p,
\qquad uv=p-2.
\]

\begin{enumerate}
\item 设三角形 \(APB\) 的重心为 \(G=(X,Y)\). 由重心坐标公式，
\[
X=\frac{p+u+v}{3}=p,
\]
且
\[
Y=\frac{p-2+u^2+v^2}{3}=\frac{p-2+(u+v)^2-2uv}{3}=\frac{4p^2-p+2}{3}.
\]
由于 \(X=p\)，消去 \(p\) 得重心 \(G\) 的轨迹方程为
\[
Y=\frac{4X^2-X+2}{3}.
\]
即 \(G\) 的轨迹方程为 \(4x^2-x+2-3y=0\).

\item 抛物线的焦点为 \(F=(0,\frac14)\). 记
\[
\bs{e}_u=\frac{\overrightarrow{FA}}{|FA|},
\qquad
\bs{e}_v=\frac{\overrightarrow{FB}}{|FB|}.
\]
由
\[
|FA|=\sqrt{u^2+\left(u^2-\frac14\right)^2}=u^2+\frac14,
\]
可得
\[
\bs{e}_u=\frac{\left(u,u^2-\frac14\right)}{u^2+\frac14},
\qquad
\bs{e}_v=\frac{\left(v,v^2-\frac14\right)}{v^2+\frac14}.
\]
令 \(S=u+v\)，\(Q=uv\)，则 \(S=2p\)，\(Q=p-2\)，即 \(2Q=S-4\). 记 \(D=(u^2+\frac14)(v^2+\frac14)>0\)，直接通分得
\[
\bs{e}_u+\bs{e}_v=\frac{\left(S\left(Q+\frac14\right),2Q^2-\frac18\right)}{D}.
\]
由 \(Q=\dfrac{S-4}{2}\) 及 \(\overrightarrow{FP}=\left(\dfrac S2,\dfrac{2S-9}{4}\right)\)，有
\[
\bs{e}_u+\bs{e}_v=\frac{2S-7}{2D}\overrightarrow{FP}.
\]
当 \(2S-7\ne0\) 时，\(FP\) 平行于两条单位向量 \(\bs{e}_u\)，\(\bs{e}_v\) 之和，故直线 \(FP\) 是两射线 \(FA\)，\(FB\) 的角平分线，因而 \(\angle PFA=\angle PFB\).

当 \(2S-7=0\) 时，\(p=\dfrac74\)，且 \(u\)，\(v\) 满足 \(4r^2-14r-1=0\). 此时 \(\overrightarrow{FP}=\left(\dfrac74,-\dfrac12\right)\)，并且对 \(r=u\) 或 \(r=v\)，均有
\[
\overrightarrow{FP}\cdot\left(r,r^2-\frac14\right)=\frac{14r-4r^2+1}{8}=0.
\]
所以 \(FP\perp FA\)，且 \(FP\perp FB\)，两个角都等于 \(\dfrac\pi2\). 综上，\(\angle PFA=\angle PFB\).
\end{enumerate}
\end{solution}
